\chapter{碳纳米管阵列自动化形貌表征与工艺—形貌关系分析}
\enchapter{Automated Morphology Characterization and Process-Morphology Analysis of Carbon Nanotube Arrays}
\label{chap:morphology_method}

\section{引言}
\ensection{Introduction}

第三章主要研究了碳纳米管阵列科研分析中的知识组织与检索增强问题，为后续工艺参数分析和结果解释提供了基础。要讨论工艺参数变化与阵列形貌演化之间的关系，还需要先获得稳定、可比较的样品级形貌表征结果。SEM 图像是观察碳纳米管阵列形貌的主要依据，但图像中普遍存在结构交叠、局部遮挡和边界模糊等情况，单纯依赖人工判读难以支持批量样品分析和定量比较。

本章的分析建立在具体实验样品、工艺参数记录和对应图像数据之上。首先说明实验样品体系、工艺记录方式和图像数据来源，并据此构建统一的样品级数据对象。

随后，本文围绕样品比较和后续工艺—形貌分析需求，建立样品级自动化表征流程。该流程以 SEM 图像为输入，通过前景分割、骨架分析和几何计算提取覆盖密度、取向度、表观直径、平均曲率、波曲度和迂曲度等关键特征，形成可用于后续分析的样品级特征表示。

在获得关键形貌表征结果后，本文结合第三章的知识检索结果开展工艺—形貌关系分析，依次进行参数筛选、统计关系分析、关系建模和结果解释。本章按实验样品与数据获取、自动化关键形貌表征和工艺—形貌关系分析三个部分展开。

\section{实验样品与数据获取}
\ensection{Samples and Data Acquisition}

\subsection{实验试剂、仪器与参数记录}
\ensubsection{Experimental Reagents, Instruments and Parameter Records}

本章实验所用主要仪器和设备见表~\ref{tab:4_1_instruments}，主要实验试剂和材料见表~\ref{tab:4_2_reagents}。

\begin{table}[htbp]
    \centering
    \caption{主要实验仪器和设备}
    \label{tab:4_1_instruments}
    \small
    \renewcommand{\arraystretch}{1.45}
    \setlength{\tabcolsep}{0pt}
    \begin{tabularx}{\textwidth}{@{}>{\centering\arraybackslash}m{0.30\textwidth}>{\centering\arraybackslash}m{0.26\textwidth}>{\centering\arraybackslash}m{0.40\textwidth}@{}}
        \toprule
        名称 & 型号 & 生产厂家 \\
        \midrule
        双靶磁控溅射仪 & VTC-600-2HD & 沈阳科晶 \\
        真空气氛管式电炉 & SK-GO8123KF & 天津中环电炉 \\
        气体质量流量控制系统 & ZL-04H & 天津中环电炉 \\
        台式场发射显微镜 & Phenom Pharos G2 & 飞纳电镜 \\
        氢气发生器 & KPS-H & 科普生 \\
        \bottomrule
    \end{tabularx}
\end{table}

\begin{table}[htbp]
    \centering
    \caption{主要实验试剂和材料}
    \label{tab:4_2_reagents}
    \small
    \renewcommand{\arraystretch}{1.45}
    \setlength{\tabcolsep}{0pt}
    \begin{tabularx}{\textwidth}{@{}>{\centering\arraybackslash}m{0.34\textwidth}>{\centering\arraybackslash}m{0.22\textwidth}>{\centering\arraybackslash}m{0.40\textwidth}@{}}
        \toprule
        名称 & 分子式/规格 & 生产厂家 \\
        \midrule
        铁靶材 & Fe & 北京晶迈中科材料技术有限公司 \\
        氧化铝靶材 & \ce{Al2O3} & 北京晶迈中科材料技术有限公司 \\
        硅单抛氧化片 & Si & 浙江立晶硅材料有限公司 \\
        石英舟 & 65 $\times$ 150 $\times$ 5 mm & 东华石英 \\
        石英管 & 1000 $\times$ 74 $\times$ 3 mm & 东华石英 \\
        高纯氩气 & Ar & 林德气体（厦门）有限公司 \\
        高纯乙烯 & \ce{C2H4} & 中山市天科气体科技有限公司 \\
        \bottomrule
    \end{tabularx}
\end{table}

本章样品涉及的主要工艺参数及其取值见表~\ref{tab:4_3_process_parameter_ranges}。

\begin{table}[htbp]
    \centering
    \caption{本章主要工艺参数及取值}
    \label{tab:4_3_process_parameter_ranges}
    \small
    \renewcommand{\arraystretch}{1.45}
    \setlength{\tabcolsep}{0pt}
    \begin{tabularx}{\textwidth}{@{}>{\centering\arraybackslash}m{0.28\textwidth}>{\centering\arraybackslash}m{0.28\textwidth}>{\centering\arraybackslash}m{0.40\textwidth}@{}}
        \toprule
        工艺阶段 & 参数 & 取值 \\ 
        \midrule
        催化层预制 & \ce{Al2O3} 功率 & 50--200 W \\
        催化层预制 & \ce{Al2O3} 厚度 & 5--30 nm \\
        催化层预制 & 铁功率 & 2--50 W \\
        催化层预制 & 铁厚度 & 0.5--2.75 nm \\
        退火预处理 & 退火温度 & 600 ℃ \\
        退火预处理 & 退火时间 & 15--60 min \\
        退火预处理/生长阶段 & Ar 流量 & 200--600 sccm \\
        退火预处理/生长阶段 & \ce{H2} 流量 & 100--300 sccm \\
        生长阶段 & \ce{C2H4} 流量 & 50--150 sccm \\
        生长阶段 & 生长温度 & 750 ℃ \\
        生长阶段 & 生长时间 & 180 min \\
        \bottomrule
    \end{tabularx}
\end{table}

\subsection{样品制备与图像数据获取流程}
\ensubsection{Sample Preparation and Image Data Acquisition Workflow}

本章样品均来自基于磁控溅射预制 \ce{Al2O3}/Fe 催化层的管式炉 CVD 生长体系。制备时，先在 Si 基底表面磁控溅射氧化铝支撑层，再沉积不同功率和厚度条件下的 Fe 催化层。随后将样片置于管式炉中心，在 Ar/\ce{H2} 气氛下完成升温和退火预处理，并在设定生长温度下引入 \ce{C2H4} 作为碳源进行生长，获得具有不同取向和弯曲状态的碳纳米管阵列样品。生长完成后，将阵列侧面竖直固定于样品台上，通过 SEM 采集侧视图像用于后续形貌分析。为保证样品级比较的可比性，后续分析主要保留基本生长窗口相对一致条件下形成的样品。

样品制备完成后，采用场发射 SEM 采集碳纳米管阵列图像，并同步记录样品编号、采集位置以及放大倍率、像素尺寸、探测模式、加速电压和工作距离等成像元数据。本研究主要采用二次电子成像模式，该模式对表面起伏和边界细节更为敏感\cite{Jackman2013,Zhang2023LVSEM}，更适合后续前景分割、骨架提取和形貌特征计算。为减小成像条件变化对结果的干扰，同一批样品尽量保持一致的观察条件。

图~\ref{fig:4_1_magnification_examples} 展示了代表性样品在不同放大倍率下的图像。较低放大倍率有利于观察阵列整体排列和宏观结构分布，但对局部边界和交叠关系的分辨能力有限。较高放大倍率能够更清楚地呈现局部边界和细节，但视场范围明显缩小，且更容易受到局部失焦和成像波动影响，不利于样品级整体统计。综合考虑整体结构信息和局部几何细节，本文统一采用 5 万倍 SEM 图像作为后续自动化表征的主要输入，相关长度和尺寸量依据图像元数据中的像素尺寸完成尺度换算。

\begin{figure}[H]
    \centering
    \includegraphics[width=0.98\textwidth]{D:/Apaper/paperformat/figures/fig4-1-magnification-grid.png}
    \caption{不同放大倍率下碳纳米管阵列侧视 SEM 图像示意。a 为 1 万倍，b 为 2 万倍，c 为 5 万倍，d 为 10 万倍。}
    \label{fig:4_1_magnification_examples}
\end{figure}

原始图像获取后，首先进行人工筛选，剔除失焦、污染较重、边界区域不完整以及存在明显充电伪影的图像，保留结构区域清晰、适合开展样品级形貌分析的图像。与此同时，对每个样品同步记录工艺参数，并建立图像编号、样品编号、工艺参数记录及成像条件之间的对应关系，保证后续结果能够回溯到具体实验样品。

图~\ref{fig:4_2_data_pipeline} 展示了本章样品级数据的构建与集成过程。实验样品、工艺参数记录和对应的 SEM 图像经过图像筛选、样品绑定和实验数据库整理后，形成后续表征与分析使用的样品级数据对象。

\begin{figure}[H]
    \centering
    \makebox[\textwidth][c]{%
    \resizebox{0.92\textwidth}{!}{%
    \begin{tikzpicture}[
        >=Stealth,
        node distance=0.85cm and 1.15cm,
        data/.style={
            draw=black!45,
            rounded corners=4pt,
            align=center,
            minimum height=1.25cm,
            text width=2.35cm,
            fill=black!3,
            font=\small
        },
        merge/.style={
            draw=black!45,
            rounded corners=4pt,
            align=center,
            minimum height=1.25cm,
            text width=2.65cm,
            fill=black!5,
            font=\small
        },
        line/.style={draw=black!60, -{Stealth[length=3mm,width=2mm]}, line width=0.7pt}
    ]
        \node[data] (sample) {实验样品};
        \node[data, right=1.15cm of sample] (params) {工艺参数记录};
        \node[data, below=1.55cm of params] (sem) {SEM 图像采集};
        \node[merge, right=1.25cm of params] (bind) {样品编号与\\图像绑定};
        \node[data, below=1.55cm of bind] (screen) {图像筛选};
        \node[data, right=1.25cm of bind] (db) {数据库集成};
        \node[merge, right=1.25cm of db] (input) {样品级数据对象\\表征与分析输入};

        \coordinate (fork) at ($(sample.east)+(0.45,0)$);

        \draw[draw=black!60, line width=0.7pt] (sample.east) -- (fork);
        \draw[line] (fork) -- (params.west);
        \draw[line] (fork) |- (sem.west);
        \draw[line] (params.east) -- (bind.west);
        \draw[line] (sem.east) -- (screen.west);
        \draw[line] (screen.north) -- (bind.south);
        \draw[line] (bind.east) -- (db.west);
        \draw[line] (db.east) -- (input.west);
    \end{tikzpicture}}
    }
    \caption{本章样品级数据构建与数据库集成流程示意图}
    \label{fig:4_2_data_pipeline}
\end{figure}

实验样品、工艺参数和图像结果由此建立一一对应关系，并整理为后续表征与分析所需的样品级数据对象。分割模型训练与评估使用相应图像子集，样品级形貌表征和工艺—形貌关系分析则基于整理后的样品图像开展。

\section{碳纳米管阵列自动化形貌表征方法}
\ensection{Automated Morphology Characterization of CNT Arrays}

\subsection{问题定义与总体流程}
\ensubsection{Problem Definition and Overall Workflow}

碳纳米管阵列 SEM 图像中常伴随细长结构交叠、局部遮挡、边界模糊和灰度分布不均等问题，样品级形貌量的稳定提取因此较为困难。围绕样品级比较与定量表征需求，本文构建关键形貌自动化表征流程。该流程先从 SEM 图像中提取稳定前景区域，再计算覆盖、取向、尺度和弯曲相关特征，最终形成样品级形貌特征表示，整体过程如图~\ref{fig:4_3_morphology_workflow} 所示。

\begin{figure}[H]
    \centering
    \makebox[\textwidth][c]{%
    \resizebox{0.94\textwidth}{!}{%
    \begin{tikzpicture}[
        node distance=0.52cm,
        >=Stealth,
        stage/.style={
            draw=black!45,
            rounded corners=4pt,
            align=center,
            minimum height=1.82cm,
            text width=3.25cm,
            fill=black!3,
            font=\small
        },
        stageC/.style={
            draw=black!45,
            rounded corners=4pt,
            align=center,
            minimum height=1.82cm,
            text width=3.45cm,
            fill=black!3,
            font=\small
        },
        arrow/.style={->, line width=0.85pt, draw=black!55}
    ]
        \node[stage] (sem) {SEM 图像输入\\阵列原始图像};
        \node[stage, right=of sem] (seg) {前景分割\\前景区域提取};
        \node[stageC, right=of seg] (feat) {形貌特征提取\\整体性与弯曲特征};
        \node[stage, right=of feat] (vec) {样品级形貌特征\\后续分析输入};

        \draw[arrow] (sem) -- (seg);
        \draw[arrow] (seg) -- (feat);
        \draw[arrow] (feat) -- (vec);
    \end{tikzpicture}}
    }
    \caption{碳纳米管阵列关键形貌自动化表征流程示意图}
    \label{fig:4_3_morphology_workflow}
\end{figure}

前景分割决定后续骨架与几何统计的输入质量。本节先给出能够较好保持细长结构的前景分割方法，再说明样品级形貌特征的提取与构建，最后通过实验比较验证整套表征方法的有效性。

\subsection{基于 CNTSegNet-clDice 的前景分割}
\ensubsection{Foreground Segmentation Based on CNTSegNet-clDice}

前景分割用于从复杂 SEM 图像中稳定提取碳纳米管区域，并为后续骨架提取和几何统计提供输入。本节关注的重点是细长结构的连通性和形态完整性，因为局部断裂、粘连或边界偏移会继续传递到后续统计过程中，降低样品级结果的稳定性。

碳纳米管在 SEM 图像中通常呈现细长、密集、交叠和缠结的二维投影，相邻个体边界往往难以稳定区分，局部区域还伴随遮挡、边界模糊和较低信噪比。在这种条件下，前景分割更适合作为自动化表征的入口任务。实例级分离不仅需要更高的标注成本，也难以在交叉和黏连区域给出稳定真值。类似的纤维网络图像分析也表明，单体分离与连续追踪本身就是独立且困难的问题\cite{Verho2025FiberSeparation,SosaRey2022OpenFiberSeg}。因此，本文先通过前景分割稳定提取可供后续骨架和几何分析使用的碳纳米管区域。

本文的分割方法由弱标签构建、分割主干和训练约束三部分组成。参考已有方向引导弱监督分割思路\cite{Nguyen2023CNTSegNet}，本文采用自适应高斯阈值生成分割监督信号，以降低精细像素级标注成本。流程为先对去除底部信息条后的主体区域进行限制对比度自适应直方图均衡化（CLAHE）增强和高斯平滑\cite{Zuiderveld1994CLAHE}，再采用自适应高斯阈值得到二值弱标签，并与主体区域图像同步裁剪为固定尺寸图块，用于模型训练。在当前实验设置下，前景分割网络训练共使用 100 张统一筛选后的 5 万倍 SEM 图像，训练所需弱标签均由上述流程自动生成。弱标签生成中的增强、平滑和自适应阈值参数在全部训练图像上保持一致，以保证监督信号构造口径统一。随后，本文在高分辨率适配的 ResNet34-UNet 主干上引入方向分布约束和 clDice 连通性约束，构建 CNTSegNet-clDice 前景分割框架，如图~\ref{fig:4_4_segmentation_framework} 所示。

\begin{figure}[H]
    \centering
    \includegraphics[width=0.98\textwidth]{D:/Apaper/paperformat/figures/fig4-4-segmentation-framework.png}
    \caption{CNTSegNet-clDice 前景分割框架示意图}
    \label{fig:4_4_segmentation_framework}
\end{figure}

前景分割主干采用高分辨率适配的 ResNet34-UNet\cite{He2016ResNet,Ronneberger2015UNet}。其中，ResNet34 通过残差连接增强深层特征提取的稳定性，能够在参数规模和表达能力之间保持较好平衡，适合当前样品量条件下的多尺度语义编码。U-Net 通过逐级上采样与跳跃连接保留浅层空间细节，更适合恢复碳纳米管这类细长目标的边界和局部形态。本文的高分辨率适配主要体现在两点。一是将首层卷积改为单通道输入并将步长调整为 1，以减弱早期下采样带来的细节损失。二是在解码末端进一步融合浅层特征，并采用轻量分割头完成输出，以加强对细长边界和局部结构的恢复。设输入 SEM 图像为 $I\in\mathbb{R}^{H\times W}$，网络输出的前景概率图记为
\begin{equation}
    P_{\theta}(I)\in[0,1]^{H\times W}
\end{equation}
其中，$\theta$ 表示网络参数。

仅靠上述主干仍难稳定保持高弯曲样品中的细长结构。前期实验表明，仅依赖方向引导时，模型在高弯曲、局部方向变化明显且结构交叠较多的区域仍易出现响应减弱、边界偏移和结构断裂。其原因在于，方向一致性先验更适用于局部主方向较稳定的区域，而高弯曲区域中的方向变化更快，局部估计也更易受到噪声、遮挡和交叠干扰，单纯依赖方向信息难以稳定保持这类区域的结构连续性。

为增强高弯曲区域中细长结构的连续性，本文引入 clDice 连通性约束。clDice 常用于细长目标分割，其做法是在软骨架空间比较预测结果与参考前景的一致性，从而约束中心线覆盖和整体连通关系\cite{Shit2021clDice}。在训练过程中，软骨架由可微软骨架算子计算得到。与仅衡量区域重叠的损失相比，clDice 同时关注预测中心线是否落在参考前景内部，以及参考中心线是否被预测结果持续覆盖。记 $\mathrm{Skel}_{s}(\cdot)$ 为可微软骨架算子，则拓扑精度 $T_{\mathrm{prec}}$ 和拓扑召回 $T_{\mathrm{sens}}$ 分别定义为
\begin{equation}
    \begin{aligned}
        T_{\mathrm{prec}}&=
        \frac{\sum_{x\in\Omega}\mathrm{Skel}_{s}\!\left(P_{\theta}(I)\right)(x)\tilde{G}(x)+\varepsilon}
        {\sum_{x\in\Omega}\mathrm{Skel}_{s}\!\left(P_{\theta}(I)\right)(x)+\varepsilon},\\
        T_{\mathrm{sens}}&=
        \frac{\sum_{x\in\Omega}\mathrm{Skel}_{s}(\tilde{G})(x)P_{\theta}(I)(x)+\varepsilon}
        {\sum_{x\in\Omega}\mathrm{Skel}_{s}(\tilde{G})(x)+\varepsilon}
    \end{aligned}
\end{equation}
据此，clDice 损失写为
\begin{equation}
    \mathcal{L}_{\mathrm{clDice}}=
    1-\frac{2T_{\mathrm{prec}}T_{\mathrm{sens}}}
    {T_{\mathrm{prec}}+T_{\mathrm{sens}}+\varepsilon}
\end{equation}
该项损失在区域重叠之外进一步约束细长结构的中心线连续性，因此更适合高弯曲区域中的结构保持。为便于比较，仅采用方向引导训练的版本记为 CNTSegNet，引入 clDice 约束的版本记为 CNTSegNet-clDice。

为兼顾基础前景定位精度与细长结构保持能力，本文采用分阶段混合损失训练策略。训练共分为两个阶段，前 6 个 epoch 仅优化 $\mathcal{L}_{\mathrm{Dice}}$，使模型先在弱标签监督下学习稳定的前景背景区分和区域定位。若在训练初期同时引入方向分布约束和 clDice 约束，早期预测中的噪声、伪响应和局部断裂会直接影响方向直方图统计与软骨架估计，不利于训练收敛。随后 3 个 epoch 以上一阶段权重为初始点继续训练，再联合引入方向分布损失与 clDice 损失，分别用于校正整体方向分布和强化高弯曲区域中的中心线覆盖与连通保持。设弱标签二值前景图为 $\tilde{G}\in\{0,1\}^{H\times W}$，图像定义域为 $\Omega$，则 Dice 损失为\cite{Milletari2016VNet}
\begin{equation}
    \mathcal{L}_{\mathrm{Dice}}=
    1-\frac{2\sum_{x\in\Omega}P_{\theta}(I)(x)\tilde{G}(x)+\varepsilon}
    {\sum_{x\in\Omega}P_{\theta}(I)(x)+\sum_{x\in\Omega}\tilde{G}(x)+\varepsilon}
\end{equation}
其中，$\varepsilon$ 为防止分母为零的微小常数。训练第一阶段的目标函数写为
\begin{equation}
    \mathcal{L}_{\mathrm{stage1}}=\mathcal{L}_{\mathrm{Dice}}
\end{equation}
第二阶段的目标函数写为
\begin{equation}
    \mathcal{L}_{\mathrm{stage2}}=
    \lambda_1\mathcal{L}_{\mathrm{Dice}}+
    \lambda_2\mathcal{L}_{\mathrm{Ori}}+
    \lambda_3\mathcal{L}_{\mathrm{clDice}}
\end{equation}
其中，$\mathcal{L}_{\mathrm{Ori}}$ 表示方向分布损失。记 $\mathcal{H}_{\mathrm{fft}}(\cdot)\in\mathbb{R}^{B}$ 为基于二维傅里叶变换功率谱统计得到的方向直方图，则
\begin{equation}
    \mathcal{L}_{\mathrm{Ori}}=
    \frac{1}{B}\sum_{b=1}^{B}
    \left(
    \mathcal{H}_{\mathrm{fft}}\!\left(P_{\theta}(I)\right)_{b}
    -
    \mathcal{H}_{\mathrm{fft}}(I)_{b}
    \right)^2
\end{equation}
其中，$B$ 表示方向直方图的分箱数。方向直方图统一在半圆角度范围内统计，并在损失计算前归一化为概率分布，以保证预测结果与原始灰度图在同一方向口径下进行比较。该项损失通过约束预测结果与原始灰度图在整体方向分布上的一致性，增强模型对阵列主方向结构的感知。本文最终取 $\lambda_1=0.6$、$\lambda_2=10^{-7}$、$\lambda_3=0.1$。这一改动不改变分割主干，而是通过调整训练目标，提高模型在复杂形貌区域中的结构保持能力。

在推理阶段，采用固定阈值 $\tau_m=0.7$ 将前景概率图二值化，得到后续形貌分析所用的二值前景图
\begin{equation}
    M(x)=\mathbb{I}\!\left(P_{\theta}(I)(x)\geq\tau_m\right)
\end{equation}
其中，$\mathbb{I}(\cdot)$ 为示性函数。本文采用统一灰度图块输入、自适应高斯阈值弱标签监督和 Adam 优化器\cite{Kingma2015Adam}，学习率设为 $5\times10^{-4}$。训练与推理的主要设置，包括分阶段混合损失训练、训练集规模、测试集评价方式和推理阈值等，汇总见表~\ref{tab:4_4_segmentation_results}。

\begin{table}[htbp]
    \centering
    \caption{前景分割模型的主要设置}
    \label{tab:4_4_segmentation_results}
    \small
    \setlength{\tabcolsep}{5pt}
    \renewcommand{\arraystretch}{1.32}
    \begin{tabular*}{\textwidth}{@{\extracolsep{\fill}}>{\centering\arraybackslash}m{3.1cm}>{\centering\arraybackslash}m{10.1cm}}
        \toprule
        项目 & 设置 \\
        \midrule
        主干网络 & ResNet34-UNet \\
        输入数据 & $768\times768$ 灰度中心图块 \\
        监督信号 & 自适应高斯阈值弱标签 \\
        优化器与学习率 & Adam，$5\times10^{-4}$ \\
        训练策略 & 分阶段混合损失训练（Dice→Dice+Ori+clDice） \\
        数据划分 & 训练集为 100 张图像，测试集固定独立划分并以 Dice 评价 \\
        推理阈值 & 0.7 \\
        \bottomrule
    \end{tabular*}
\end{table}

\subsection{样品级形貌特征提取与构建}
\ensubsection{Sample-Level Morphology Feature Extraction and Construction}

在获得稳定的碳纳米管二值前景图后，本文提取样品级形貌特征。覆盖密度、取向度和表观直径反映样品的覆盖状态、方向有序性和尺度信息，平均曲率、波曲度和迂曲度用于描述整体弯曲形貌。为进行路径级几何统计，本文先对二值前景进行骨架化，再在骨架基础上完成拓扑清理和有效分支提取。图~\ref{fig:4_5_feature_pipeline} 给出了弯曲相关特征提取中的关键中间结果。

\begin{figure}[H]
    \centering
    \includegraphics[width=0.95\textwidth]{D:/Apaper/paperformat/figures/fig4-5-feature-pipeline-grid.png}
    \caption{弯曲相关特征提取关键中间结果图。a 为原始 SEM 图像，b 为预测前景图，c 为原始骨架，d 为拓扑清理中，e 为清理后骨架，f 为连接点邻域排除。}
    \label{fig:4_5_feature_pipeline}
\end{figure}

\noindent　　\textbf{（1）原始骨架生成与拓扑清理}\par

本文首先对二值前景进行骨架化，将碳纳米管区域转换为单像素宽的中心线网络，以保留主要连通关系并用于后续分支级几何统计\cite{VanDerWalt2014ScikitImage}。设二值前景图为 $M$，其原始骨架记为
\begin{equation}
    B_{\mathrm{raw}}=\mathrm{Skel}(M)
\end{equation}
原始骨架通常仍包含短小分量、端部毛刺以及连接点附近的局部伪分支，因此本文进一步进行拓扑清理，以获得用于弯曲相关特征计算的稳定有效分支。该过程首先删除很短、很小且与主拓扑弱相关的孤立骨架分量，以降低局部噪声和边界伪影的影响。由于本文所用图像放大倍率统一，拓扑清理中的尺度阈值均依据图像元数据中的像素尺寸换算到像素空间。设像素尺寸为 $s_{\mathrm{px}}$，单位为 \si{nm/pixel}。本文以 15 nm 作为拓扑清理中的经验管径尺度，并据此定义期望管径像素尺度
\begin{equation}
    d_{\mathrm{tube,expected}}^{(\mathrm{px})}=
    \max\!\left(1.5,\;\frac{15}{s_{\mathrm{px}}}\right)
\end{equation}
进一步地，最小路径长度阈值和最小像素点数阈值分别设为
\[
    L_{\min}^{(\mathrm{px})}=\max\!\left(3.0,\;2.0\,d_{\mathrm{tube,expected}}^{(\mathrm{px})}\right),\qquad
    N_{\min}=\max\!\left(4,\;\mathrm{round}\!\left(1.2\,d_{\mathrm{tube,expected}}^{(\mathrm{px})}\right)\right)
\]
对任一 8 邻域连通骨架分量，若其像素点数小于 $N_{\min}$，或路径长度小于 $L_{\min}^{(\mathrm{px})}$，则将其判定为弱相关小分量并予以删除。

从端点沿骨架追踪至最近连接点，若端点到连接点的路径长度满足 $L_{\mathrm{spur}}<K\cdot d_{\mathrm{tube,expected}}^{(\mathrm{px})}$，则将其判定为端部毛刺并予以删除。式中，$K$ 为剪枝系数，本文默认取 $K=3$。上述尺度阈值依据当前样品在 5 万倍视野下的常见表观管径量级及预实验结果确定，用于抑制明显短小伪分支并保留主要连通主干。完成端部毛刺剪枝后，再进行一次细化，以消除连接点附近可能残留的局部厚边或团块状结构，得到拓扑更稳定的清理后骨架。随后，在清理后骨架上依次进行连接点邻域排除、分支有序化和最小长度过滤，提取用于弯曲相关特征计算的有效分支集合 $\mathcal{B}^{*}$\cite{NunezIglesias2018Skan}。由于交叉和黏连区域中的跨交点延拓本身具有歧义，本文在排除连接点邻域后仅对可用分支进行局部几何分析\cite{Verho2025FiberSeparation,SosaRey2022OpenFiberSeg}。为评估连接点邻域排除后仍保留的信息量，定义有效骨架保留率为
\begin{equation}
    \eta_{\mathrm{keep}}=
    \frac{\sum_{b_k\in\mathcal{B}^{*}}L_k}{L_{\mathrm{raw}}+\varepsilon}
\end{equation}
其中，$L_{\mathrm{raw}}$ 表示原始骨架总长度。该统计量用于衡量拓扑清理后仍参与几何分析的骨架长度占比，从而避免复杂区域被过度剔除。

除用于构建核心形貌特征向量的 6 项特征外，本文还在拓扑清理后的骨架上统计骨架连接点数、骨架总长度和连接点密度，用于辅助评估前景结果的结构连通复杂度与骨架保持情况。记清理后骨架中的连接点集合为 $\mathcal{J}$，则有
\[
    N_{\mathrm{conn}} = |\mathcal{J}|,\qquad
    L_{\mathrm{skel}} = \sum_{b_k\in\mathcal{B}^{*}}L_k,\qquad
    D_{\mathrm{conn}} = \frac{100\,N_{\mathrm{conn}}}{L_{\mathrm{skel}}+\varepsilon}
\]
其中，$L_{\mathrm{skel}}$ 表示骨架总长度，单位为 $\mu$m。$D_{\mathrm{conn}}$ 表示每 100 $\mu$m 骨架长度对应的连接点数。这三项统计量主要服务于后续分割结果比较与结果核查，不纳入工艺—形貌关系分析所用的核心特征向量。

\noindent　　\textbf{（2）覆盖密度与取向度计算}\par

覆盖密度和取向度沿原始几何路径计算，不以拓扑清理后的骨架结果替代原始几何表示。在覆盖密度计算中，记前景像素面积为 $A_{\mathrm{fg}}$，图像主体区域面积为 $A_{\mathrm{ROI}}$，则覆盖密度定义为
\begin{equation}
    \rho=\frac{A_{\mathrm{fg}}}{A_{\mathrm{ROI}}}\times 100\%
\end{equation}
在取向度计算中，本文基于原始骨架 $B_{\mathrm{raw}}$ 上的局部主方向样本估计 Herman 取向因子（HOF）\cite{Brandley2018,Yadav2024}。样品主方向由当前图像主体区域的骨架方向分布自适应确定。为降低图像整体旋转带来的方向歧义，计算时同时考察互相正交的两个候选参考方向，并选取取向度更稳定的一组结果。设局部方向与样品主方向之间的夹角为 $\phi$，则整体取向度记为
\begin{equation}
    \mathrm{HOF}=\frac{3\langle \cos^2\phi\rangle-1}{2}
\end{equation}
该处理可在样品存在轻微旋转时提高阵列主方向识别的稳定性。

\noindent　　\textbf{（3）表观直径估计}\par

表观直径同样沿原始几何路径估计，以避免拓扑剪枝对局部尺度统计造成额外偏差。本文通过距离变换读取原始骨架点处的半径，并在闭运算与骨架化后对半径样本进行四分位距过滤，以降低边界毛刺和局部异常点的影响。设过滤后的像素半径集合为 $r_{\mathrm{px}}^{\mathrm{filt}}$，图像元数据中的像素尺寸为 $s_{\mathrm{px}}$，单位为 \si{nm/pixel}，则样品的表观直径估计为
\begin{equation}
    d_{\mathrm{app}}=2\,\mathrm{median}\!\left(r_{\mathrm{px}}^{\mathrm{filt}}\right)\,s_{\mathrm{px}}
\end{equation}
其中，$d_{\mathrm{app}}$ 的单位为 nm。

\noindent　　\textbf{（4）弯曲相关特征统计}\par

弯曲相关特征沿拓扑清理路径进行统计。本文采用平均曲率、波曲度和迂曲度从不同层面描述阵列弯曲形貌。其中，平均曲率反映局部几何转折强度，波曲度反映分支整体起伏幅度与波长之间的相对关系，迂曲度反映路径长度相对于端点直线距离的绕行程度。三者分别对应局部弯折、整体波状起伏和路径曲折程度，因此能够形成互补表征。

对于每条有效分支，本文先进行有序化与轻度平滑，再在离散采样点上利用三点法估计局部曲率。设第 $k$ 条有效分支记为 $b_k=\{p_{k,1},p_{k,2},\ldots,p_{k,n_k}\}$，对于其中第 $i$ 个内部采样点，定义
\[
    \mathbf{ab}_{k,i}=p_{k,i}-p_{k,i-1},\qquad
    \mathbf{bc}_{k,i}=p_{k,i+1}-p_{k,i},\qquad
    \mathbf{ca}_{k,i}=p_{k,i+1}-p_{k,i-1}
\]
则局部离散曲率近似为
\begin{equation}
    \kappa_{k,i}=
    \frac{2\left|
    \mathbf{ab}_{k,i}^{x}\mathbf{bc}_{k,i}^{y}-
    \mathbf{ab}_{k,i}^{y}\mathbf{bc}_{k,i}^{x}
    \right|}
    {\left\|\mathbf{ab}_{k,i}\right\|_2
    \left\|\mathbf{bc}_{k,i}\right\|_2
    \left\|\mathbf{ca}_{k,i}\right\|_2+\varepsilon}
\end{equation}
其中，$\varepsilon$ 为防止分母为零的微小常数。考虑到数字曲线上的曲率估计对噪声和方向扰动较为敏感\cite{Kerautret2009Curvature,Vieilleville2009Tangent,Frette2009Curvature}，本文在局部曲率计算前先对分支路径进行轻度平滑，再以分支内局部曲率的平均值表征该分支的平均弯曲程度如下。
\begin{equation}
    \kappa_{k}^{\mathrm{br}}=
    \frac{1}{n_k-2}\sum_{i=2}^{n_k-1}\kappa_{k,i}
\end{equation}
设第 $k$ 条分支的路径长度为 $L_k$，则样品级有效平均曲率定义为
\begin{equation}
    \kappa_{\mathrm{eff}}=
    \frac{\sum_{b_k\in\mathcal{B}^{*}}\sqrt{L_k}\,\kappa_{k}^{\mathrm{br}}}
    {\sum_{b_k\in\mathcal{B}^{*}}\sqrt{L_k}}
\end{equation}
样品级平均曲率在统计后依据图像元数据中的像素尺寸换算为 $\mu m^{-1}$。样品级聚合采用 $\sqrt{L_k}$ 而非等权或 $L_k$ 线性加权，以避免短碎分支权重过高，同时减弱超长分支对结果的主导作用。

由于实际分支的横向起伏通常并不满足严格的正弦形态，本文不采用显式正弦拟合，而是依据分支主轴方向上的等效波高和等效波长估计波曲度。对于单条有效分支，先对其坐标执行主成分分析，确定主轴方向和法向方向，并分别投影得到纵向坐标 $s$ 和横向偏移 $d(s)$。对 $d(s)$ 进行线性去趋势和轻度平滑后，检测交替出现的峰谷结构，并将相邻峰谷在主轴方向上的距离记为等效波长 $\lambda$，对应的横向偏移差记为等效波高 $h$。于是，单个波段的波曲度定义为
\begin{equation}
    w=\frac{h}{\lambda}
\end{equation}
对同一分支内的多个波段，按各波段对应的 $\lambda$ 加权平均得到分支级波曲度。对样品中的多条分支，再按分支长度的平方根加权，得到样品级波曲度
\begin{equation}
    w_k=\frac{\sum_{j}\lambda_{k,j}\,w_{k,j}}{\sum_{j}\lambda_{k,j}},\qquad
    w_{\mathrm{eff}}=
    \frac{\sum_{b_k\in\mathcal{B}^{*}}\sqrt{L_k}\,w_k}
    {\sum_{b_k\in\mathcal{B}^{*}}\sqrt{L_k}}
\end{equation}
迂曲度则采用路径长度与端点弦长之比进行度量，用于反映整条分支相对于端点直线距离的绕行程度，定义如下。
\begin{equation}
    \tau_k=\frac{L_k}{\left\|p_{k,n_k}-p_{k,1}\right\|_2+\varepsilon},\qquad
    \tau_{\mathrm{eff}}=
    \frac{\sum_{b_k\in\mathcal{B}^{*}}\sqrt{L_k}\,\tau_k}
    {\sum_{b_k\in\mathcal{B}^{*}}\sqrt{L_k}}
\end{equation}
样品级波曲度和迂曲度的多分支聚合同样按分支长度的平方根加权，以与曲率统计保持一致，并兼顾短分支与长分支之间的权重平衡。

\noindent　　\textbf{（5）样品级特征构建}\par

据此，样品级形貌特征向量可写为
\[
    \mathbf{z}=
    \left[\rho,\,\mathrm{HOF},\,d_{\mathrm{app}},\,\kappa_{\mathrm{eff}},\,w_{\mathrm{eff}},\,\tau_{\mathrm{eff}}\right]
\]
其中，前三项来自原始几何路径，后三项来自拓扑清理路径。该表示同时保留了阵列整体有序性、覆盖状态、尺度信息和样品级弯曲形貌。表~\ref{tab:4_5_feature_vector_summary} 汇总了各特征的组成及其含义。

\begin{table}[htbp]
    \centering
    \caption{样品级形貌特征的组成及其含义}
    \label{tab:4_5_feature_vector_summary}
    \small
    \setlength{\tabcolsep}{2pt}
    \begin{tabularx}{\textwidth}{@{}>{\centering\arraybackslash}m{2.0cm}>{\centering\arraybackslash}m{4.8cm}>{\centering\arraybackslash}m{2.0cm}>{\centering\arraybackslash}m{1.5cm}>{\centering\arraybackslash}m{3.7cm}@{}}
        \toprule
        特征 & 物理含义 & 计算路径 & 单位 & 后续用途 \\
        \midrule
        覆盖密度 $\rho$ & 表征前景区域覆盖程度 & 原始路径 & \% & 反映覆盖状态与生长均匀性 \\
        取向度 HOF & 表征阵列主方向一致程度 & 原始路径 & 无量纲 & 反映阵列整体有序性 \\
        表观直径 $d_{\mathrm{app}}$ & 表征图像尺度下的等效宽度 & 原始路径 & nm & 反映局部尺度特征 \\
        有效平均曲率 $\kappa_{\mathrm{eff}}$ & 表征局部几何转折强度 & 清理路径 & $\mu m^{-1}$ & 描述样品级弯曲强度 \\
        波曲度 $w_{\mathrm{eff}}$ & 表征整体起伏幅度与波长的相对关系 & 清理路径 & 无量纲 & 区分平直与波曲形貌 \\
        迂曲度 $\tau_{\mathrm{eff}}$ & 表征路径相对端点直线距离的绕行程度 & 清理路径 & 无量纲 & 描述卷曲与缠绕程度 \\
        \bottomrule
    \end{tabularx}
\end{table}

\subsection{实验结果与分析}
\ensubsection{Experimental Results and Analysis}

前两节给出了前景分割方法和形貌特征提取流程。本节统一比较弱标签参考、CNTSegNet 和 CNTSegNet-clDice 三类前景结果，依次分析结构统计差异、弯曲特征与人工分级的一致性以及处理效率。这里的 Dice 系数和交并比（Intersection over Union, IoU）衡量的是模型预测结果与弱标签参考结果之间的区域重叠一致性，而非相对于人工精细标注真值的绝对分割精度，因此主要用于比较不同模型在同一弱标签监督口径下的相对表现。

先比较三类前景结果在高弯曲样品中的结构保持差异。图~\ref{fig:4_6_segmentation_comparison} 给出了四组代表样本在原图、弱标签、CNTSegNet 结果和 CNTSegNet-clDice 结果下的前景输入对比。

\begin{figure}[H]
    \centering
    \includegraphics[width=\textwidth]{D:/Apaper/paperformat/figures/fig4-6-segmentation-comparison-grid.png}
    \caption{高弯曲代表样本的前景输入对比图。a--d 分别对应原图、弱标签、CNTSegNet 结果和 CNTSegNet-clDice 结果，各列从左至右对应四组代表样本。}
    \label{fig:4_6_segmentation_comparison}
\end{figure}

对于高弯曲样品，前景分割不仅要区分前景与背景，还要尽量保持弯折结构的连续性和局部拓扑稳定。后续骨架提取对局部断裂、细丝粘连和边界偏移较为敏感，一旦弯折段在分割阶段被破坏，误差便会继续传递到分支提取和弯曲特征计算中。局部断裂会缩短可用路径，细丝粘连会改变分支走向，连接点附近的伪连接还会干扰波长、路径长度和端点距离等量的估计。从图中可以看到，在局部交叠较多、弯折更明显的代表样本中，CNTSegNet-clDice 相比 CNTSegNet 更能保持细长结构的连续形态。图中的示例主要展示高弯曲难例中的局部差异，整体结果见表~\ref{tab:4_6_segmentation_structure_compare}。

\begin{table}[htbp]
\centering
\caption{三类分割结果在测试图像集上的平均一致性与骨架结构统计对比}
\label{tab:4_6_segmentation_structure_compare}
\footnotesize
\setlength{\tabcolsep}{2.6pt}
\renewcommand{\arraystretch}{1.28}
\begin{tabular*}{\textwidth}{@{\extracolsep{\fill}}>{\centering\arraybackslash}m{2.95cm}S[table-format=1.3]S[table-format=1.3]S[table-format=2.2]S[table-format=1.3]S[table-format=5.2]S[table-format=4.2]S[table-format=3.2]}
\toprule
结果来源 & {Dice} & {IoU} & {覆盖密度/\%} & {取向度} & {\makecell[c]{平均骨架\\连接点数}} & {\makecell[c]{连接点密度\\(个/100$\mu$m)}} & {\makecell[c]{平均骨架\\总长度/$\mu$m}} \\
\midrule
弱标签参考 & {\textemdash} & {\textemdash} & 52.83 & 0.791 & 10888.26 & 2734.29 & 375.71 \\
CNTSegNet & 0.891 & 0.804 & 44.19 & 0.874 & 4555.86 & 1604.04 & 276.33 \\
CNTSegNet-clDice & 0.897 & 0.813 & 44.84 & 0.867 & 5074.16 & 1713.76 & 286.91 \\
\bottomrule
\end{tabular*}
\end{table}

表中各项数值均为对测试图像集逐图统计后取平均值得到。弱标签参考不参与 Dice 和 IoU 计算，仅保留其结构统计，作为统一比较基准。CNTSegNet 和 CNTSegNet-clDice 则同时比较区域重叠一致性和骨架结构统计。相较于 CNTSegNet，CNTSegNet-clDice 的 Dice 由 0.891 提升至 0.897，IoU 由 0.804 提升至 0.813，增幅都不大。这说明在同一弱标签监督口径下，两种模型在区域层面的重叠差异有限，clDice 的主要作用不在于显著改变前景面积，而在于改善细长结构的保持方式。

这一差异在骨架结构统计中更明显。引入 clDice 后，覆盖密度由 44.19\% 提升至 44.84\%，平均骨架总长度由 276.33 $\mu$m 提升至 286.91 $\mu$m，平均骨架连接点数由 4555.86 提升至 5074.16，连接点密度由 1604.04 增加至 1713.76。区域重叠指标只有小幅变化，而骨架长度、连接点数和连接点密度同时增长，说明新增保留下来的前景主要集中在原本容易断裂的弯折段和连接区域附近。若只是增加零散伪前景或简单加厚边界，骨架总长度可能上升，但连接点密度未必会同步增加，因此这组变化更符合连通性约束改善细长主干保持的预期。

与弱标签参考相比，其覆盖密度达到 52.83\%，明显高于两种学习模型，平均骨架总长度、平均骨架连接点数和连接点密度也整体更高。这与弱标签前景保留更宽有关。自适应阈值作用于增强后的灰度图时，会将部分低对比细丝、局部黏连和边界扩张区域一并纳入前景，使骨架网络显得更稠密，相应结构统计量也随之偏大。因此，弱标签在这些指标上的高值更多反映前景保留尺度较宽，不能直接等同于结构质量更优。本文更关心的是在统一弱监督口径下模型对有效连通结构的保留能力。就这一点而言，CNTSegNet-clDice 在保持 Dice 和 IoU 基本稳定的同时，恢复了更多可用于后续骨架分析的连续结构。其取向度由 0.874 略降至 0.867，也说明模型在整体方向集中度略有降低的同时，保留了更多局部弯折和连接结构。对于碳纳米管这类细长连通目标，仅依赖区域重叠指标难以充分反映分割结果对后续形貌分析的支持程度，骨架结构统计更能揭示模型是否保留了与弯曲特征计算相关的有效几何信息。

在比较前景结果后，进一步考察三类结果是否会影响后续弯曲特征与人工分级的一致性。一致性验证所用数据集共包含 150 张样品图像，其中平直型、波曲型和卷曲型样本分别为 40、80 和 30 张。这 150 张图像用于样品级弯曲特征与人工分级一致性验证，不对应前景分割网络的训练集规模。前景分割网络训练见 4.3.2 节，训练集为 100 张图像。三类标签用于描述样品级整体弯曲状态。平直型主要对应阵列整体较为顺直、局部扰动较弱的样品。波曲型对应存在较明显波状起伏、但尚未表现出强烈缠绕的中等弯曲样品。卷曲型则对应整体路径绕行更明显、局部缠结和高弯曲区域更突出的样品。图~\ref{fig:4_7_manual_grading_examples} 给出了三类样本的代表性形貌。

标注工作由 3 名具有相关研究背景的领域人员完成，采用双人交叉标注与第三人复核相结合的方式。人工等级标注面向样品整体弯曲形貌，主要依据整张图像所呈现的总体弯曲水平和整体结构表现进行判断。

\begin{figure}[H]
    \centering
    \includegraphics[width=\textwidth]{D:/Apaper/paperformat/figures/fig4-6-manual-grading-examples.png}
    \caption{人工弯曲等级划分示意图。a 为平直型，b 为波曲型，c 为卷曲型。}
    \label{fig:4_7_manual_grading_examples}
\end{figure}

在这一部分，仍统一比较弱标签参考、CNTSegNet 和 CNTSegNet-clDice，并对曲率、波曲度和迂曲度三个单指标分别进行阈值分级验证。为减少阈值选择与性能评估使用同一批样本所带来的乐观偏差，本文采用两折交叉方式开展一致性验证。按类别分层将样本固定划分为两个子集，每次选取其中 75 张作为训练集，根据训练集人工分级结果确定三等级划分中的两条阈值边界，再将该阈值固定应用于另外 75 张测试图像，完成自动分级与人工分级的一致性统计，随后交换训练集与测试集重复一次，从而保证全部 150 张图像均在独立测试条件下参与评估。考虑到三类样本数量并不均衡，结果采用准确率（Accuracy）、宏平均 F1 值（Macro-F1）、平衡准确率（Balanced Accuracy）和加权 Kappa 系数（Weighted Kappa）四项指标进行评价\cite{Cohen1968WeightedKappa}。下文先比较三类结果下的混淆矩阵，再对表现最优的 CNTSegNet-clDice 结果进行定量汇总。

\begin{figure}[H]
    \centering
    \includegraphics[width=0.98\textwidth]{D:/Apaper/paperformat/figures/fig4-8-confusion-matrix-grid.png}
    \caption{不同前景获取方式与弯曲特征组合下的人工分级一致性混淆矩阵。列从左至右依次为曲率、波曲度和迂曲度，行从上至下依次为 CNTSegNet-clDice、CNTSegNet 和弱标签参考。}
    \label{fig:4_8_confusion_matrices}
\end{figure}

图~\ref{fig:4_8_confusion_matrices} 展示了三类前景输入与三种弯曲特征组合下的人工分级一致性混淆矩阵。整体来看，九组结果几乎都没有出现平直型与卷曲型之间的大规模直接错判，误差主要集中在相邻等级之间，而且错分样本大多向中间等级波曲型聚集。这说明三等级划分的主要难点不在于极端形貌的识别，而在于过渡形貌边界的稳定确定。波曲型在定义上处于平直型与卷曲型之间，因此更容易吸收两侧样本，这也解释了混淆矩阵中第二列普遍较宽的现象。

按行比较三类前景输入，可以看到更好的连通性保持主要改善了波曲型与卷曲型之间的区分。以迂曲度为例，卷曲型的正确识别率由弱标签参考的 60.0\% 提升至 CNTSegNet 的 66.7\%，并进一步提升至 CNTSegNet-clDice 的 73.3\%。波曲型的正确识别率也由 90.0\% 提升至 92.5\% 和 96.2\%。这说明减少局部断裂后，基于路径长度与端点距离的整体路径统计更加稳定。波曲度也呈现类似趋势，其中波曲型的正确识别率由 80.0\% 提升至 86.2\% 和 92.5\%，说明连通性改进有助于稳定波段结构识别。相比之下，曲率对前景输入变化的收益较弱，平直型与波曲型之间仍存在较明显混淆，即使在 CNTSegNet-clDice 结果下，平直型仍有 45.0\% 被划入波曲型，说明该指标更易受局部细小弯折、骨架扰动和曲线平滑细节影响。

按列比较三个弯曲特征，迂曲度的对角线最集中，且对卷曲型样本的保持最好，更适合承担整图排序与主判别任务。波曲度在中间等级上的稳定性较好，更适合补充表征波状起伏。曲率则更接近局部弯折强度指标，更适合作为辅助描述。

\begin{table}[htbp]
    \centering
    \caption{CNTSegNet-clDice 结果下不同弯曲特征与人工分级结果的一致性比较}
    \label{tab:4_7_manual_auto_comparison}
    \small
    \setlength{\tabcolsep}{5pt}
    \begin{tabularx}{\textwidth}{@{}>{\centering\arraybackslash}m{4.2cm}*{4}{>{\centering\arraybackslash}m{2.0cm}}@{}}
        \toprule
        指标 & Accuracy & Macro-F1 & Balanced Accuracy & Weighted Kappa \\
        \midrule
        曲率 & 0.7467 & 0.7239 & 0.7042 & 0.6955 \\
        波曲度 & 0.8133 & 0.7957 & 0.7667 & 0.7732 \\
        迂曲度 & 0.9133 & 0.9008 & 0.8819 & 0.9006 \\
        \bottomrule
    \end{tabularx}
\end{table}

表~\ref{tab:4_7_manual_auto_comparison} 从整体指标上验证了混淆矩阵中的观察。考虑到三类样本数量并不均衡，Balanced Accuracy 和 Macro-F1 更能反映各等级识别是否均衡，Weighted Kappa 则用于衡量三等级排序的一致性。迂曲度在四项指标上均最高，Accuracy 达到 0.9133，Balanced Accuracy 为 0.8819，Weighted Kappa 为 0.9006，说明其优势并非仅来自波曲型样本数量较多，而是在三类样本上都保持了较稳定的区分能力。波曲度的 Accuracy 和 Macro-F1 分别为 0.8133 和 0.7957，整体低于迂曲度，但在中间等级上的识别仍较稳定，适合补充描述波状起伏。曲率的四项指标最低，尤其 Balanced Accuracy 仅为 0.7042，这与其在混淆矩阵中更易将平直型和波曲型混淆的现象一致。由此可见，弯曲相关特征可以作为样品级弯曲表征量，其中迂曲度更适合作为主判据，波曲度和曲率更适合作为补充描述。

最后比较人工粗筛与自动分级的处理效率和结果一致性。本文以表现最稳定的迂曲度分级结果作为自动方法代表，并与人工粗筛结果进行对比。该自动分级准确率对应前述 CNTSegNet-clDice 前景输入下迂曲度单指标分级结果的 Accuracy。表~\ref{tab:4_8_efficiency_comparison} 汇总了两种方法在同一批 150 张样品图像上的总耗时、平均耗时和准确率，其中准确率均以人工三级标注结果作为参考标签。自动耗时按单张图像顺序执行自动化表征与基于迂曲度的分级流程统计，不计人工复核与结果整理时间。结果显示，人工粗筛总耗时为 1960\,s，平均约 13.07\,s/张，准确率为 67.3\%。基于迂曲度的自动分级总耗时为 1548\,s，平均约 10.32\,s/张，准确率为 91.33\%。在当前实现条件下，自动方法既保持了较高的一致性，也在批量样品处理时表现出一定的效率优势。

\begin{table}[htbp]
    \centering
    \caption{人工粗筛与自动分级的效率及准确率比较}
    \label{tab:4_8_efficiency_comparison}
    \small
    \setlength{\tabcolsep}{5pt}
    \begin{tabularx}{\textwidth}{@{}>{\centering\arraybackslash}m{3.2cm}*{3}{>{\centering\arraybackslash}m{3.2cm}}@{}}
        \toprule
        方法 & 总耗时（s） & 平均耗时/张（s） & 准确率 \\
        \midrule
        人工粗筛 & 1960 & 13.07 & 67.3\% \\
        自动分级（迂曲度） & 1548 & 10.32 & 91.33\% \\
        \bottomrule
    \end{tabularx}
\end{table}

\section{知识增强支持下的工艺—形貌关系分析}
\ensection{Knowledge-Enhanced Process-Morphology Analysis}

前一节已经完成样品级形貌指标的提取与一致性验证。本节在当前 \ce{Al2O3}/Fe 热 CVD 样品体系内，围绕取向度、平均曲率、波曲度和迂曲度等关键形貌指标，讨论工艺参数与形貌特征之间的对应关系。考虑到原始工艺记录涉及催化层预制、退火和生长等多个环节，后续讨论先对工艺变量范围作适度收敛，再开展统计分析和关系建模。本节同时沿用第三章构建的工艺分析任务模板及其检索结果，用于说明分析范围收敛的依据，并对基于样品级实验数据得到的统计分析和关系建模结果提供机理解释支持。

\subsection{工艺参数筛选与分析范围收敛}
\ensubsection{Process Parameter Screening and Scope Convergence}

原始记录同时包含 \ce{Al2O3} 参数、Fe 参数、气氛条件、温度和时间等多类条件。若将这些变量同时展开，样品可比性和讨论主线都会受到影响。因此，这里先从全变量层面观察整体相关格局，再据此确定后续统计分析所保留的核心变量。

\begin{figure}[H]
    \centering
    \includegraphics[width=0.80\textwidth]{D:/Apaper/paperformat/figures/fig4-6-all-process-heatmap.png}
    \caption{全部候选工艺变量与关键形貌特征的相关性热图}
    \label{fig:4_9_all_process_heatmap}
\end{figure}

图~\ref{fig:4_9_all_process_heatmap} 给出了全部候选工艺变量与样品级形貌特征之间的相关结构。可以看到，\ce{Al2O3} 功率/厚度、气氛档位以及多数背景生长条件与几项形貌特征的相关性整体较弱。相比之下，退火时间与取向度、波曲度和迂曲度的相关性更突出，铁功率虽与弯曲相关特征保持弱正相关，但强度明显低于退火时间，铁厚度的单变量相关性则更弱。这说明在全变量范围内，退火时间更接近当前样品集中形貌变化方向的主要对应变量，铁功率更适合作为退火时间之外的补充调节变量，铁厚度则更多提供催化剂负载水平的背景差异。后续分析仍保留这三项变量，但在结果解释上，将以退火时间和铁功率作为更适合展开讨论的主线。这里的收敛并不意味着其他条件不重要，而是意味着它们在当前样品体系中更适合作为背景生长窗口和样品可比性的控制条件。

\definecolor{kgInputBg}{RGB}{244,247,252}
\definecolor{kgInputTag}{RGB}{98,117,147}
\definecolor{kgEvidenceBg}{RGB}{248,245,240}
\definecolor{kgEvidenceTag}{RGB}{141,112,84}
\definecolor{kgSummaryBg}{RGB}{241,248,245}
\definecolor{kgSummaryTag}{RGB}{85,125,106}
\definecolor{kgResultBg}{RGB}{227,240,233}
\newcommand{\kgfield}[2]{%
    {%
        \renewcommand{\arraystretch}{1.10}%
        \setlength{\tabcolsep}{4pt}%
        \noindent\begin{tabularx}{\linewidth}{@{}>{\raggedright\arraybackslash}p{0.14\linewidth}X@{}}
            \textbf{#1} & #2 \\
        \end{tabularx}\par\smallskip
    }%
}
\newcommand{\kgsectiontitle}[3]{%
    \setlength{\fboxsep}{6pt}%
    \noindent\colorbox{#1}{%
        \parbox{\dimexpr\linewidth-2\fboxsep\relax}{%
            {\color{#2}\bfseries #3}%
        }%
    }\par\smallskip
}
\newcommand{\kgsectionbox}[4]{%
    \setlength{\fboxsep}{7pt}%
    \noindent\fcolorbox{black!18}{white}{%
        \parbox{\dimexpr\linewidth-2\fboxsep-2\fboxrule\relax}{%
            \kgsectiontitle{#1}{#2}{#3}%
            #4
        }%
    }\par
}
\newcommand{\kgcompactfields}[3]{%
    {%
        \renewcommand{\arraystretch}{1.12}%
        \setlength{\tabcolsep}{4pt}%
        \noindent\begin{tabularx}{\linewidth}{@{}>{\raggedright\arraybackslash}p{0.20\linewidth}X@{}}
            \textbf{任务类型} & #1 \\
            \textbf{自然语言输入} & #2 \\
            \textbf{查询结构化约束} & #3 \\
        \end{tabularx}\par
    }%
}
\newcommand{\kgchaincard}[3]{%
    {%
        \renewcommand{\arraystretch}{1.10}%
        \setlength{\tabcolsep}{4pt}%
        \noindent\begin{tabularx}{\linewidth}{@{}>{\raggedright\arraybackslash}p{0.13\linewidth}X@{}}
            \textbf{#1} & #2 \\
            {\color{black!68}\textbf{文献依据}} & {\color{black!68}#3} \\
        \end{tabularx}\par
    }%
    {\color{black!12}\rule{\linewidth}{0.4pt}}\par\smallskip
}
\newcommand{\kgresult}[1]{%
    \setlength{\fboxsep}{6pt}%
    \noindent\colorbox{kgResultBg}{%
        \parbox{\dimexpr\linewidth-2\fboxsep\relax}{%
\textbf{归纳结果} #1%
        }%
    }\par
}
\newcommand{\kgpanel}[3]{%
    \setlength{\fboxsep}{8pt}%
    \noindent\fcolorbox{black!12}{#2}{%
        \parbox{\dimexpr\linewidth-2\fboxsep-2\fboxrule\relax}{%
            {\color{#3}\bfseries #1}\par\medskip
        }%
    }%
}
\newcommand{\kgpanelcontent}[2]{%
    \setlength{\fboxsep}{8pt}%
    \noindent\fcolorbox{black!12}{#1}{%
        \parbox{\dimexpr\linewidth-2\fboxsep-2\fboxrule\relax}{#2}%
    }\par
}
\begin{figure}[H]
\centering
\footnotesize
\setlength{\fboxsep}{6pt}
\begin{minipage}{0.97\linewidth}
\kgsectionbox{kgInputBg}{kgInputTag}{输入}{
\kgcompactfields{工艺—形貌分析中的变量范围收敛。}{在当前基于 \ce{Al2O3}/Fe 催化层的热 CVD 碳纳米管阵列样品体系中，原始工艺记录同时包含催化层预制、退火预处理和生长阶段的多项条件。若目标是在保持背景生长窗口基本一致的前提下，分析工艺参数与样品级形貌特征之间的对应关系，哪些变量更适合作为后续统计分析和关系建模的保留变量。}{$E_q=\{$\ce{Al2O3} 功率/厚度，铁功率/厚度，Ar，\ce{H2}，\ce{C2H4}，退火起始温度，生长温度，生长时间，退火时间$\}$；$C_q=\{$磁控溅射，\ce{Al2O3}/Fe 催化层，管式炉热 CVD$\}$；$Y_q=\{$取向度，平均曲率，波曲度，迂曲度，样品可比性$\}$。}
}
\vspace{0.4em}
\kgsectionbox{kgEvidenceBg}{kgEvidenceTag}{TCCER 证据链输出}{
    \kgchaincard{证据链 1}{生长温度、碳源供给和总体生长时间首先决定阵列生长窗口、产率和整体生长行为。}{垂直排列碳纳米管阵列的生长机理研究与综述普遍将温度、碳源和生长时间视为基础过程条件\cite{Hata2004,Liu2021EcoMat,Wu2025VACNTReview}。}
    \kgchaincard{证据链 2}{铁厚度会影响催化层退火后的颗粒尺寸、分布和密度基础，并进一步影响阵列微结构与生长表现。}{已有研究表明，Fe 催化层厚度会影响碳纳米管阵列的微结构、密度及相关生长表现\cite{Yang2016ALDThicknessCN,Zhang2017BufferLayer}。}
    \kgchaincard{证据链 3}{退火时间对应催化层预处理阶段的还原、脱润湿和颗粒演化窗口，进而影响后续成核与阵列生长表现。}{已有研究表明，预退火过程会改变 Fe 催化颗粒形成状态，并进一步影响阵列生长结果\cite{Jung2013PreAnnealingFeCNT,Lee2019CatalystDewetting}。}
    \kgchaincard{证据链 4}{铁功率会影响前驱 Fe 薄膜的结构、表面形貌和均匀性，从而可能继续影响后续颗粒化状态。}{磁控溅射研究表明，溅射功率会显著影响 Fe 薄膜的结构与表面形貌\cite{Aleksanyan2023RFMagnetron,Zhou2020FeSputterPower}。}
}
\vspace{0.4em}
\kgsectionbox{kgSummaryBg}{kgSummaryTag}{范围收敛说明}{
\kgfield{归纳依据}{结合上述证据，在当前样品体系中，与后续形貌响应联系更直接的是催化层形成与演化过程。退火时间、铁厚度和铁功率都进入这一过程链条，而 \ce{Al2O3} 参数、气氛条件、生长温度和总体生长时间在本文中已作为背景生长窗口固定，用于保证样品可比性。}
\kgresult{因此，后续工艺—形貌分析围绕退火时间、铁功率和铁厚度展开。该图用于说明变量范围收敛的知识依据，不直接承担变量主次判断。}
}
\end{minipage}
\caption{变量范围收敛的知识辅助说明}
\label{fig:param_convergence}
\end{figure}

图~\ref{fig:param_convergence} 给出了第三章任务模板返回的证据链整理结果。结合这些证据可以看到，退火时间、铁功率和铁厚度都与催化层形成及其后续形貌响应联系更直接，因此有必要保留进入后续分析。相比之下，\ce{Al2O3} 参数、Ar/\ce{H2}/\ce{C2H4} 条件、生长温度和总体生长时间在当前样品体系中更多起背景控制作用。变量主次判断主要由后续统计和建模结果给出。

\subsection{工艺参数与关键形貌特征的统计关系分析}
\ensubsection{Statistical Analysis of Process Parameters and Key Morphology Features}

在上述变量范围下，本文在固定 \ce{Al2O3} 功率与厚度、气氛条件、退火温度、生长温度和生长时间的前提下，围绕铁功率、铁厚度和退火时间构建了 204 个样品级分析对象，对应 21 组铁功率—铁厚度组合。每个样品对应若干张经过筛选的可用 5 万倍 SEM 图像。为尽量降低 SEM 视野位置差异对形貌统计的影响，图像筛选时优先保留拍摄位置相近且视野范围一致的可用图像。本文先对单张图像提取形貌特征，再对同一样品内筛选后的可用图像结果取算术平均，得到该样品的样品级形貌表征值。后续相关性分析和关系建模均以该样品级平均值为基础。统一采用 5 万倍图像，是为了保证形貌指标在像素尺度和视野范围上的可比性。分析输入为铁功率、铁厚度和退火时间，输出关注取向度、平均曲率、波曲度和迂曲度。统计分析按三个层次展开，即整体相关、退火时间分组趋势和现有铁功率—铁厚度组合的描述性响应分布。由于铁厚度在当前样品中主要集中于两个主层，其余厚度多为补充点，因此该组合图更适合作为对响应区间的描述性观察，而不宜直接理解为完整二维响应面。

首先采用 Spearman 秩相关系数\cite{Spearman1904Association}，对工艺参数与关键形貌特征之间的整体相关结构进行统计分析。设两组变量的秩序列分别为 $R_i$ 和 $S_i$，则其相关系数定义为
\begin{equation}
    \rho_s=
    \frac{\sum_{i=1}^{n}(R_i-\bar{R})(S_i-\bar{S})}
    {\sqrt{\sum_{i=1}^{n}(R_i-\bar{R})^2}\sqrt{\sum_{i=1}^{n}(S_i-\bar{S})^2}}
\end{equation}
其中，$\bar{R}$ 和 $\bar{S}$ 分别表示两组秩序列的均值。图~\ref{fig:4_10_spearman_correlation_heatmap} 给出了收敛变量与关键形貌特征之间的相关性热图。可以看到，退火时间与取向度、平均曲率、波曲度和迂曲度的 Spearman 系数分别为 $-0.54$、$0.51$、$0.54$ 和 $0.58$，说明退火时间增加时，阵列整体取向下降而弯曲相关特征增强，这一对应关系在当前样品集上最为稳定。取向度与平均曲率、波曲度和迂曲度均呈负相关，其中与波曲度和迂曲度的相关性更强，表明有序性降低通常伴随弯曲增强。波曲度与迂曲度高度正相关，说明两者对弯曲变化的响应基本一致。相比之下，铁功率与形貌特征的单变量相关性整体较弱，但对弯曲相关特征仍保持弱正相关，说明其更可能作为退火时间之外的调节因素影响不同工艺窗口中的响应区间。铁厚度的单变量相关性更弱，且当前样品在该维度上的覆盖主要集中于两个主层，因此更适合作为后续解释中的补充背景变量。由此可见，退火时间更适合用来刻画形貌变化方向，而铁功率更适合作为进一步解释响应差异的补充变量。

\begin{figure}[H]
    \centering
    \includegraphics[width=0.74\textwidth]{D:/Apaper/paperformat/figures/fig4-7-converged-heatmap.png}
    \caption{收敛变量与关键形貌特征的相关性热图}
    \label{fig:4_10_spearman_correlation_heatmap}
\end{figure}

\begin{figure}[H]
    \centering
    \includegraphics[width=0.76\textwidth]{D:/Apaper/paperformat/figures/fig4-8-anneal-grid.png}
    \caption{退火时间对关键形貌指标的影响趋势。a 为取向度，b 为平均曲率，c 为波曲度，d 为迂曲度。}
    \label{fig:4_11_anneal_time_trend}
\end{figure}

按退火时间分组后，这一趋势更为直观。图~\ref{fig:4_11_anneal_time_trend} 中，随着退火时间由 0.25 h 延长至 1.00 h，取向度整体下降，波曲度和迂曲度整体上升，说明退火过程推动阵列由相对有序状态向更明显的弯曲状态演化。平均曲率在 30 min 之后整体抬升，但波动更大，进一步说明它对局部扰动更敏感。因而在样品级尺度上，波曲度和迂曲度更适合作为弯曲状态的主要表征量。

退火时间之外，铁功率是更值得继续展开的补充变量。它在单变量层面不表现为稳定的单调趋势，但不同功率组对应的形貌分布区间仍存在差异。相比之下，铁厚度在当前样品中的覆盖主要集中于 1.0 nm 和 2.0 nm 两个主层，其余厚度多为补充点，因此更适合作为背景变量保留，而不宜继续作为与退火时间同等展开的主轴。

\begin{figure}[H]
    \centering
    \includegraphics[width=0.84\textwidth]{D:/Apaper/paperformat/figures/fig4-12-power-boxplot-grid.png}
    \caption{铁功率对关键形貌指标的分组分布。a 为取向度，b 为平均曲率，c 为波曲度，d 为迂曲度。}
    \label{fig:4_12_response_grid}
\end{figure}

图~\ref{fig:4_12_response_grid} 按铁功率分组给出了四项关键形貌指标的样品分布。可以看到，不同功率组之间的中位数和四分位区间并不完全重合，说明铁功率虽然不是主导形貌变化方向的一阶变量，但仍会改变不同样品组的响应区间。对于取向度，5 W、30 W 和 70 W 组整体较高，而 2 W 和 40 W 组分布更低且离散更明显。对于平均曲率、波曲度和迂曲度，中高功率组的整体分布普遍高于低功率组，但组间仍存在重叠。这说明铁功率更适合作为补充调节变量理解，其作用主要体现为改变样品分布位置和离散程度，而不是单独决定弯曲增强方向。

综合相关热图、退火时间分组结果和铁功率分组箱线图，可以将三项变量的作用区分为两个层次。退火时间主要对应取向下降和弯曲增强的总体方向，铁功率主要改变不同功率组的分布位置和离散程度，铁厚度则更多提供催化剂负载水平的背景差异。基于这一结果，后续建模重点检验上述变量角色能否在非线性关系中保持稳定。

\subsection{工艺—形貌关系建模与解释}
\ensubsection{Process-Morphology Modeling and Explanation}

前面的统计结果给出了变量作用方向，但还不能说明这些工艺条件在当前样品子集中能否稳定对应到样品级形貌指标。碳纳米管阵列的工艺参数存在耦合，弯曲特征还容易受到局部几何扰动和成像噪声影响。因此，本文引入随机森林回归模型\cite{Breiman2001RandomForests}，检验铁功率、铁厚度和退火时间在当前样品子集中能否对取向度、平均曲率、波曲度和迂曲度形成可学习对应关系。由于输入已统一为样品级平均特征，模型评估采用样品级 5 折交叉验证，即每次用其中 4 折训练、1 折验证，并轮流完成全部样品的预测。模型性能以交叉验证下的决定系数（$R^2$）和均方根误差（RMSE）衡量，同时结合 SHAP 结果比较各变量对不同形貌特征预测的贡献，其主要超参数设置见表~\ref{tab:4_9_rf_settings}。

\begin{table}[htbp]
    \centering
    \caption{随机森林回归模型的主要超参数设置}
    \label{tab:4_9_rf_settings}
    \small
    \renewcommand{\arraystretch}{1.45}
    \setlength{\tabcolsep}{0pt}
    \begin{tabularx}{\textwidth}{@{}>{\centering\arraybackslash}m{0.30\textwidth}>{\centering\arraybackslash}m{0.66\textwidth}@{}}
        \toprule
        参数类别 & 参数设置 \\
        \midrule
        模型 & 随机森林回归 \\
        决策树数量 & 300 \\
        随机种子 & 42 \\
        最大深度 & 不限制 \\
        节点最小分裂样本数 & 2 \\
        叶节点最小样本数 & 1 \\
        最大特征数 & 1.0 \\
        自助采样 & 是 \\
        分裂准则 & 平方误差 \\
        \bottomrule
    \end{tabularx}
\end{table}

为便于从各形貌特征角度比较预测效果和变量贡献，下面按取向度、平均曲率、波曲度和迂曲度分别给出建模结果。各组图中，(a) 为按测试样本序号排列的实测值与预测值对比，(b) 为交叉验证下的实测值与预测值散点，(c) 为 SHAP 值分布，(d) 为平均绝对 SHAP 值排序。

\begin{figure}[H]
    \centering
    \includegraphics[width=0.80\textwidth]{D:/Apaper/paperformat/figures/fig4-10-alignment-combo.png}
    \caption{取向度的建模结果与 SHAP 分析。a 为序列对比，b 为回归散点，c 为 SHAP 分布，d 为特征贡献排序。}
    \label{fig:4_13_alignment_combo}
\end{figure}

\begin{figure}[H]
    \centering
    \includegraphics[width=0.80\textwidth]{D:/Apaper/paperformat/figures/fig4-10-curvature-combo.png}
    \caption{平均曲率的建模结果与 SHAP 分析。a 为序列对比，b 为回归散点，c 为 SHAP 分布，d 为特征贡献排序。}
    \label{fig:4_14_curvature_combo}
\end{figure}

\begin{figure}[H]
    \centering
    \includegraphics[width=0.80\textwidth]{D:/Apaper/paperformat/figures/fig4-11-waviness-combo.png}
    \caption{波曲度的建模结果与 SHAP 分析。a 为序列对比，b 为回归散点，c 为 SHAP 分布，d 为特征贡献排序。}
    \label{fig:4_15_waviness_combo}
\end{figure}

\begin{figure}[H]
    \centering
    \includegraphics[width=0.80\textwidth]{D:/Apaper/paperformat/figures/fig4-11-tortuosity-combo.png}
    \caption{迂曲度的建模结果与 SHAP 分析。a 为序列对比，b 为回归散点，c 为 SHAP 分布，d 为特征贡献排序。}
    \label{fig:4_16_tortuosity_combo}
\end{figure}

取向度的预测效果最好，$R^2$ 为 $0.6225 \pm 0.1432$，RMSE 为 $0.0501 \pm 0.0079$。图~\ref{fig:4_13_alignment_combo} 显示，预测曲线能够较好跟踪样本序列上的整体起伏，散点分布也最接近对角线，说明工艺条件对阵列整体有序性的影响最容易被模型稳定捕捉。SHAP 结果进一步表明，退火时间的平均绝对 SHAP 值为 0.046，高于铁功率和铁厚度，且高退火时间主要对应负向贡献，这与 4.4.2 节中退火时间增加时取向度下降的统计结果一致。

平均曲率的 $R^2$ 为 $0.5453 \pm 0.0689$，RMSE 为 $1.1446 \pm 0.1184$。从图~\ref{fig:4_14_curvature_combo} 可以看出，平均曲率仍表现出一定可学习性，但其样本序列对比和散点分布都比取向度更分散。SHAP 排序中，退火时间的平均绝对 SHAP 值达到 1.082，明显高于铁厚度和铁功率，说明退火时间仍是解释平均曲率变化的首要变量。与此同时，高退火时间更多对应正向贡献，表明在当前样品集上，退火过程延长往往伴随局部弯折强度上升。需要说明的是，平均曲率在回归任务中的拟合程度略高于波曲度和迂曲度，但结合前述相关性和分组趋势，它在样品级统计上的波动仍然较大，更适合作为辅助性的连续弯折指标理解。

波曲度的 $R^2$ 为 $0.4799 \pm 0.0467$，RMSE 为 $0.0135 \pm 0.0013$。图~\ref{fig:4_15_waviness_combo} 显示，模型能够捕捉其整体变化方向，但离散程度已经明显高于取向度。SHAP 结果中，退火时间的平均绝对 SHAP 值为 0.010，仍居首位，铁功率和铁厚度分别为 0.005 和 0.003。高退火时间主要对应正向 SHAP 值，说明退火时间增加时，波曲度更容易升高。这与前面退火时间分组趋势图中的观察一致，也说明铁功率和铁厚度更多是在不同工艺区间内改变波曲响应的强弱和位置。

迂曲度的 $R^2$ 为 $0.4929 \pm 0.0691$，RMSE 为 $0.0059 \pm 0.0003$。图~\ref{fig:4_16_tortuosity_combo} 表明，其回归结果与波曲度接近，但散点仍保持较清楚的正相关趋势。SHAP 排序显示，退火时间的平均绝对 SHAP 值为 0.006，仍高于铁功率和铁厚度，且高退火时间同样主要对应正向贡献。这表明退火时间延长时，路径绕行程度更容易增加。结合 4.3 节中迂曲度在人工分级一致性验证中的优势，可以认为该指标不仅能稳定反映弯曲增强方向，也更适合作为样品级弯曲状态的重点指标。

四个特征合并来看，退火时间在每个回归任务中的预测贡献都最高，铁功率始终保持次一级且稳定的非零贡献，铁厚度则相对更弱。前者主要对应取向下降和弯曲增强的总体方向，铁功率更多表现为对不同工艺窗口响应强弱和位置的调节，铁厚度则更接近催化剂负载水平的背景变量。这一角色分工与第三章返回的证据链一致。相关知识证据只用于解释，不参与统计计算、模型训练或性能评估。表~\ref{tab:4_10_variable_role_summary} 汇总了统计结果、建模结果和知识证据。

\begin{table}[htbp]
    \centering
    \caption{关键工艺变量作用的统计、建模与知识证据汇总}
    \label{tab:4_10_variable_role_summary}
    \small
    \setlength{\tabcolsep}{4pt}
    \renewcommand{\arraystretch}{1.30}
    \begin{tabularx}{\textwidth}{@{}>{\centering\arraybackslash}m{2.1cm}*{4}{>{\raggedright\arraybackslash}X}@{}}
        \toprule
        变量 & 统计现象 & 建模证据 & 知识证据 & 最终角色判断 \\
        \midrule
        铁功率 & 单变量连续趋势较弱，但在弯曲相关特征上保持弱正相关，不同功率区间对应的响应位置存在差异。 & 在四个特征中的平均绝对 SHAP 值始终低于退火时间，但保持稳定非零贡献，且通常高于铁厚度。 & 铁功率影响沉积速率、膜层均匀性和铁前驱薄层形成质量，进而影响后续颗粒化状态。 & 适合作为退火时间之外的调节变量保留。 \\
        铁厚度 & 单变量相关性更弱，且当前样品在该维度上的覆盖主要集中于两个主层，其余厚度多为补充点。 & 在四个特征中的平均绝对 SHAP 值均为非零，但整体低于铁功率和退火时间。 & 铁厚度更接近铁催化薄膜的名义厚度，影响退火后颗粒尺寸、分布和密度基础。 & 适合作为催化剂负载水平的补充背景变量保留。 \\
        退火时间 & 与取向度下降呈负相关，与波曲度和迂曲度增强呈正相关，分组趋势最稳定。 & 在四个形貌特征中的平均绝对 SHAP 值均最高，交叉验证下对取向度及弯曲相关特征都保持了稳定的趋势预测。 & 退火阶段对应催化层还原、脱润湿和颗粒演化窗口，进而影响成核、生长一致性和结构稳定性。 & 适合作为解释取向变化和弯曲增强趋势的重要变量。 \\
        \bottomrule
    \end{tabularx}
\end{table}

铁功率和铁厚度都作用于催化层形成阶段，但在当前样品中的角色并不对称。铁功率更直接对应磁控溅射阶段的沉积速率、膜层均匀性和铁前驱薄层成膜质量，因此更适合作为退火时间之外的调节变量理解。铁厚度更接近名义膜厚和催化剂负载水平，影响退火后颗粒尺寸、分布和密度基础，但在当前样品中的覆盖不如铁功率充分，因此更适合作为补充背景变量保留。

退火时间的角色更直接。统计分析和建模都表明，它与取向下降、弯曲增强之间存在更稳定的对应关系。结合知识证据，退火阶段本质上是催化层由前驱体薄膜向活性颗粒转变的过程，时间变化会改变颗粒形成与演化窗口，进而影响后续阵列生长的一致性和结构稳定性。因此，在当前样品体系中，退火时间更适合作为解释取向变化和弯曲增强趋势的重要变量。

综上，在当前样品体系且背景生长窗口基本一致的前提下，铁功率、铁厚度和退火时间构成了后续工艺—形貌分析的核心变量集。其中，退火时间主要对应取向下降和弯曲增强的变化方向，铁功率主要调节不同退火窗口中的响应强弱和位置，铁厚度则更多体现催化剂负载水平的背景差异。样品级自动化表征提供了可比较的形貌指标，知识检索结果则补充了变量收敛和机制解释所需依据。

\section{小结}
\ensection{Summary}

本章首先给出了碳纳米管阵列样品和图像数据来源，并建立了面向样品级比较的自动化形貌表征流程。在当前样品体系内，铁功率、铁厚度和退火时间与取向度、平均曲率、波曲度和迂曲度之间表现出可学习的对应关系，其中退火时间与取向下降及弯曲增强之间的趋势最为稳定，铁功率主要影响不同退火窗口中的响应强弱和位置，铁厚度更多体现催化剂负载水平的背景差异。知识检索结果为变量收敛和结果解释补充了依据，由此形成了面向当前样品体系的工艺—形貌关系分析结果。
